\documentclass[11pt,a4paper]{article}
\usepackage[margin=27mm]{geometry}
\usepackage{amsmath,amssymb,amsthm}
\usepackage{unicode-math}
\setmainfont{STIXGeneral}
\setmathfont{STIX Two Math}
\usepackage[hidelinks]{hyperref}
\newtheorem{theorem}{Theorem}
\title{Regular polytopes with projective general linear groups}
\author{Human Compression project}
\date{}
\begin{document}
\maketitle

The stronger theorem for string markings with two proper ends in [U]
makes the main classification of [LS] a direct consequence.

\begin{theorem}[Leemans--Schulte]
A regular abstract polytope of rank at least four with full
automorphism group $\operatorname{PGL}_2(q)$ is the 4-simplex, and
$q=5$. Its distinguished marking consists of the adjacent transpositions
in $S_5$, of type $\{3,3,3\}$.
\end{theorem}

\begin{proof}
The intersection property makes both end groups proper: a generator
belonging to the subgroup on the remaining generators would violate
the intersection with its own subgroup of order two. By the
projective restriction and finite classification theorems of [U],
$q$ is odd, the rank is four, and the marked group is one of
$\operatorname{PSL}_2(11)$, $\operatorname{PSL}_2(19)$ or $S_5$.

For odd $q$, determinant modulo squares gives
$\operatorname{PGL}_2(q)$ a quotient of order two. The two simple
groups are therefore impossible. Thus the group is $S_5$;
$|\operatorname{PGL}_2(q)|=q(q^2-1)=120$ forces $q=5$ by strict
monotonicity. The finite classification and geometry theorems of [U]
identify its ordered generators with $(12),(23),(34),(45)$ and
establish the intersection property and uniqueness. The standard
group--polytope correspondence now gives the 4-simplex. Conversely
this marking realizes that simplex, with full group
$S_5\cong\operatorname{PGL}_2(5)$.
\end{proof}

This is [LS, Theorem 10] and its formulation for every rank
at least four. The seventeen-case semilinear census, chiral discussion
and Conjecture 11 in that paper are not among the results proved here.
The present proof uses the common paper [U] directly; it does not
assume the separate PSL application. The complete common proof and
finite program are charged once within this topic, with the
classical subgroup results and polytope correspondences stated explicitly.
The supporting Lean certificates do not constitute a formal proof of
this whole classification.

\begin{thebibliography}{9}
\bibitem[U]{U} \emph{Common theory for projective string markings},
companion paper, projective restriction, finite classification and
geometry theorems.
\bibitem[LS]{LS} D. Leemans and E. Schulte,
\emph{Polytopes with groups of type $\operatorname{PGL}_2(q)$},
arXiv:0909.1991v1, the bound nine-page author edition,
Theorem 10 on page 7.
\end{thebibliography}
\end{document}
